% GENERATED FILE. Edit canonical lessons, then run npm run generate:books.
% canonical-source-hash: fnv1a64:cda250abe3c49788
% canonical-lessons: JA-C224-gakusei, JA-C224-kyoushitsu, JA-C224-jugyou, JA-C224-oshieru, JA-C224-narau

\chapter{Students and Classrooms}
\label{ch:ja-students-and-classrooms}
\hlchaptermodality{224}

\begin{chapteropening}
\textbf{By the end of this chapter:} \emph{I can talk about students, classrooms and lessons.}

Complete the last lesson of chapter 224: I can talk about students, classrooms and lessons.
\end{chapteropening}

\section[\texorpdfstring{\ja{がくせい} (gakusei)}{がくせい (gakusei)}]{\ja{がくせい} (gakusei) — a student; a university student}

\label{lesson:JA-C224-gakusei}



\begin{quote}
Before the new one: say the Japanese for open for business, then the Japanese for not yet open.
\end{quote}

\subsection*{You'll want to know: \ja{がくせい}}
\textbf{\ja{がくせい}} — \emph{gakusei} — \textquotedblleft{}a student; a university student\textquotedblright{}.

\begin{quote}
\textbf{\ja{がくせいは} \ja{まいにち} \ja{がっこうに} \ja{いきます。}} — \emph{gakusei wa mainichi gakkō ni ikimasu} — \textquotedblleft{}The students go to school every day.\textquotedblright{}
\end{quote}

\begin{cousinweb}
\textbf{\ja{がくせい}} (\emph{gakusei}) is Sino-Japanese (kango): two Chinese characters read with their on'yomi, readings borrowed from Chinese. They are \emph{gaku}, \textquotedblleft{}learning\textquotedblright{} (Middle Chinese \emph{haewk}) and \emph{sei}, \textquotedblleft{}life, birth\textquotedblright{} (\emph{srjaeng}) (KANJIDIC2; Middle Chinese in Baxter's transcription).
\end{cousinweb}

\begin{grammarlens}[title={What you've built}]
One more word for everyday life at A2.
\end{grammarlens}

\subsection*{Your turn}
\begin{itemize}
\raggedright
  \item \emph{Say it:} \emph{gakusei}
  \item \emph{Say it:} \emph{gakusei}, once more
  \item \emph{Say it:} \emph{gakusei wa mainichi gakkō ni ikimasu} — The students go to school every day.
  \item \emph{Recall:} say the Japanese for open for business, then the Japanese for not yet open, then say \emph{gakusei} again
\end{itemize}

\begin{tcolorbox}[breakable,colback=teal!4,colframe=teal!35!black,title={Before you move on}]
What does \ja{がくせい} mean? (\textquotedblleft{}A student; a university student\textquotedblright{}.) Say it once more.
\end{tcolorbox}
\section[\texorpdfstring{\ja{きょうしつ} (kyōshitsu)}{きょうしつ (kyōshitsu)}]{\ja{きょうしつ} (kyōshitsu) — a classroom}

\label{lesson:JA-C224-kyoushitsu}



\begin{quote}
Before the new one: say the Japanese for not yet open, then the Japanese for a student.
\end{quote}

\subsection*{You'll want to know: \ja{きょうしつ}}
\textbf{\ja{きょうしつ}} — \emph{kyōshitsu} — \textquotedblleft{}a classroom\textquotedblright{}.

\begin{quote}
\textbf{\ja{がくせいは} \ja{きょうしつで} \ja{べんきょうします。}} — \emph{gakusei wa kyōshitsu de benkyō shimasu} — \textquotedblleft{}The students study in the classroom.\textquotedblright{}
\end{quote}

\begin{cousinweb}
\textbf{\ja{きょうしつ}} (\emph{kyōshitsu}) is Sino-Japanese (kango): two Chinese characters read with their on'yomi, readings borrowed from Chinese. They are \emph{kyō}, \textquotedblleft{}teach\textquotedblright{} (Middle Chinese \emph{kaewH}) and \emph{shitsu}, \textquotedblleft{}room\textquotedblright{} (\emph{syit}) (KANJIDIC2; Middle Chinese in Baxter's transcription).
\end{cousinweb}

\begin{grammarlens}[title={What you've built}]
One more word for everyday life at A2.
\end{grammarlens}

\subsection*{Your turn}
\begin{itemize}
\raggedright
  \item \emph{Say it:} \emph{kyōshitsu}
  \item \emph{Say it:} \emph{kyōshitsu}, once more
  \item \emph{Say it:} \emph{gakusei wa kyōshitsu de benkyō shimasu} — The students study in the classroom.
  \item \emph{Recall:} say the Japanese for not yet open, then the Japanese for a student, then say \emph{kyōshitsu} again
\end{itemize}

\begin{tcolorbox}[breakable,colback=teal!4,colframe=teal!35!black,title={Before you move on}]
What does \ja{きょうしつ} mean? (\textquotedblleft{}A classroom\textquotedblright{}.) Say it once more.
\end{tcolorbox}
\section[\texorpdfstring{\ja{じゅぎょう} (jugyō)}{じゅぎょう (jugyō)}]{\ja{じゅぎょう} (jugyō) — a class; a lesson}

\label{lesson:JA-C224-jugyou}



\begin{quote}
Before the new one: say the Japanese for a student, then the Japanese for a classroom.
\end{quote}

\subsection*{You'll want to know: \ja{じゅぎょう}}
\textbf{\ja{じゅぎょう}} — \emph{jugyō} — \textquotedblleft{}a class; a lesson\textquotedblright{}.

\begin{quote}
\textbf{\ja{ごごは} \ja{きょうしつで} \ja{じゅぎょうが} \ja{あります。}} — \emph{gogo wa kyōshitsu de jugyō ga arimasu} — \textquotedblleft{}In the afternoon there is a class in the classroom.\textquotedblright{}
\end{quote}

\begin{cousinweb}
\textbf{\ja{じゅぎょう}} (\emph{jugyō}) is Sino-Japanese (kango): two Chinese characters read with their on'yomi, readings borrowed from Chinese. They are \emph{ju}, \textquotedblleft{}impart\textquotedblright{} (Middle Chinese \emph{dzyuwH}) and \emph{gyō}, \textquotedblleft{}work, vocation\textquotedblright{} (\emph{ngjaep}) (KANJIDIC2; Middle Chinese in Baxter's transcription).
\end{cousinweb}

\begin{grammarlens}[title={What you've built}]
One more word for everyday life at A2.
\end{grammarlens}

\subsection*{Your turn}
\begin{itemize}
\raggedright
  \item \emph{Say it:} \emph{jugyō}
  \item \emph{Say it:} \emph{jugyō}, once more
  \item \emph{Say it:} \emph{gogo wa kyōshitsu de jugyō ga arimasu} — In the afternoon there is a class in the classroom.
  \item \emph{Recall:} say the Japanese for a student, then the Japanese for a classroom, then say \emph{jugyō} again
\end{itemize}

\begin{tcolorbox}[breakable,colback=teal!4,colframe=teal!35!black,title={Before you move on}]
What does \ja{じゅぎょう} mean? (\textquotedblleft{}A class; a lesson\textquotedblright{}.) Say it once more.
\end{tcolorbox}
\section[\texorpdfstring{\ja{おしえる} (oshieru)}{おしえる (oshieru)}]{\ja{おしえる} (oshieru) — to teach; to tell}

\label{lesson:JA-C224-oshieru}



\begin{quote}
Before the new one: say the Japanese for a classroom, then the Japanese for a class.
\end{quote}

\subsection*{You'll want to know: \ja{おしえる}}
\textbf{\ja{おしえる}} — \emph{oshieru} — \textquotedblleft{}to teach; to tell\textquotedblright{}.

\begin{quote}
\textbf{\ja{せんせいは} \ja{にほんごを} \ja{おしえます。}} — \emph{sensei wa nihongo o oshiemasu} — \textquotedblleft{}The teacher teaches Japanese.\textquotedblright{}
\end{quote}

Polite \textbf{\ja{おしえます}}, \emph{oshiemasu}; an ichidan verb.

\begin{cousinweb}
\textbf{\ja{おしえる}} (\emph{oshieru}) is a native Japanese word (wago), inherited rather than borrowed from Chinese; its kanji was attached later as a kun'yomi (JMdict, KANJIDIC2).
\end{cousinweb}

\begin{grammarlens}[title={What you've built}]
One more word for everyday life at A2.
\end{grammarlens}

\subsection*{Your turn}
\begin{itemize}
\raggedright
  \item \emph{Say it:} \emph{oshieru}
  \item \emph{Say it:} \emph{oshieru}, once more
  \item \emph{Say it:} \emph{sensei wa nihongo o oshiemasu} — The teacher teaches Japanese.
  \item \emph{Recall:} say the Japanese for a classroom, then the Japanese for a class, then say \emph{oshieru} again
\end{itemize}

\begin{tcolorbox}[breakable,colback=teal!4,colframe=teal!35!black,title={Before you move on}]
What does \ja{おしえる} mean? (\textquotedblleft{}To teach; to tell\textquotedblright{}.) Say it once more.
\end{tcolorbox}
\section[\texorpdfstring{\ja{ならう} (narau)}{ならう (narau)}]{\ja{ならう} (narau) — to take lessons; to learn from a teacher}

\label{lesson:JA-C224-narau}



\begin{quote}
Before the new one: say the Japanese for a class, then the Japanese for to teach.
\end{quote}

\subsection*{You'll want to know: \ja{ならう}}
\textbf{\ja{ならう}} — \emph{narau} — \textquotedblleft{}to take lessons; to learn from a teacher\textquotedblright{}.

\begin{quote}
\textbf{\ja{いもうとは} \ja{せんせいに} \ja{うたを} \ja{ならいます。}} — \emph{imōto wa sensei ni uta o naraimasu} — \textquotedblleft{}My younger sister takes singing lessons from a teacher.\textquotedblright{}
\end{quote}

Polite \textbf{\ja{ならいます}}, \emph{naraimasu}; a godan verb.

\begin{cousinweb}
\textbf{\ja{ならう}} (\emph{narau}) is a native Japanese word (wago), inherited rather than borrowed from Chinese; its kanji was attached later as a kun'yomi (JMdict, KANJIDIC2).
\end{cousinweb}

\begin{grammarlens}[title={What you've built}]
One more word for everyday life at A2.
\end{grammarlens}

\subsection*{Your turn}
\begin{itemize}
\raggedright
  \item \emph{Say it:} \emph{narau}
  \item \emph{Say it:} \emph{narau}, once more
  \item \emph{Say it:} \emph{imōto wa sensei ni uta o naraimasu} — My younger sister takes singing lessons from a teacher.
  \item \emph{Recall:} say the Japanese for a class, then the Japanese for to teach, then say \emph{narau} again
\end{itemize}

\begin{tcolorbox}[breakable,colback=teal!4,colframe=teal!35!black,title={Before you move on}]
What does \ja{ならう} mean? (\textquotedblleft{}To take lessons; to learn from a teacher\textquotedblright{}.) Say it once more.
\end{tcolorbox}
